\documentclass[11pt]{article}

\usepackage[margin=1.1in]{geometry}
\usepackage{amsmath,amssymb,amsthm}
\usepackage{stmaryrd}
\usepackage{mathtools}
\usepackage{tikz}
\usetikzlibrary{arrows.meta,positioning,fit,backgrounds,patterns,calc,decorations.pathmorphing}
\usepackage{listings}
\usepackage{xcolor}
\usepackage{booktabs}
\usepackage{array}
\usepackage[colorlinks=true,linkcolor=black,citecolor=black,urlcolor=blue]{hyperref}
\usepackage[expansion=false]{microtype}

\theoremstyle{plain}
\newtheorem{theorem}{Theorem}
\newtheorem{proposition}[theorem]{Proposition}
\newtheorem{lemma}[theorem]{Lemma}
\newtheorem{corollary}[theorem]{Corollary}
\newtheorem{conjecture}[theorem]{Conjecture}
\theoremstyle{definition}
\newtheorem{definition}[theorem]{Definition}
\newtheorem{remark}[theorem]{Remark}
\newtheorem{observation}[theorem]{Observation}
\newtheorem{question}{Question}

% --- typography: the rho calculus is NEVER written with the Greek letter ---
\newcommand{\rhoc}{\textsf{rho}}
\newcommand{\Chn}{\mathrm{Chn}}

\lstdefinestyle{rho}{
  basicstyle=\ttfamily\small,
  keywordstyle=\bfseries,
  morekeywords={for,new,in,def,where,contract,match,if,else,Nil},
  columns=fullflexible,
  showstringspaces=false,
  breaklines=true,
  frame=none,
  xleftmargin=1.5em,
  comment=[l]{//},
  commentstyle=\itshape\color{black!55},
}
\lstset{style=rho}

\title{\bfseries The Tube\\[.35em]
\large Topological consequences of heterotrophy in the mortal scientist framework}
\author{L.\ G.\ Meredith\\ \small F1R3FLY.io}
\date{September 8, 2026 \\ \small revised September 11, 2026}

\begin{document}
\maketitle

\begin{abstract}
A heterotroph eats in order to reach the token supply held inside another agent. Time spent
\emph{exclusively} on extracting that supply is time not spent foraging, mating, or forming and
testing hypotheses, and so carries an opportunity cost that selection can be expected to attack. The
attack available is concurrency: split the harvest into a fast acquisition phase and a slow internal
phase, and let the second run in parallel with everything else the agent does. The internal phase
leaves a residue --- the part that was not energy --- and the residue must leave. Material in, refuse
out. This note argues that the heterotroph is therefore, topologically, a \emph{tube}, that the
autotroph has no corresponding pressure, and that this is exactly the arrangement observed in the
physical world: animals are tubes, plants are not. The argument is set inside the mortal scientist
framework, where the pieces it requires --- trophic access profiles, interposed media, porosity and
partial namespace closure --- are already present. Nothing new is assumed. The claim is that the tube
is derivable, and in particular that it is a second corollary of the same proposition that already
yields ``intelligence is a heterotroph's expense''. The note is a research note in the strict sense:
it records an argument, a list of statements to attempt, and the empirical checks each would face.
The September 11 revision adds two things. The first is a refinement of the criterion: the
discriminating variable is not the rate of exchange but an occupancy ratio together with a
separability, which is what distinguishes eating from breathing and explains why respiratory systems
themselves are sometimes tubes. The second is a set of counterexamples proposed by M.\ Stay ---
bag-gutted animals, and heterotrophs outside the animals --- which are worked through one by one. One
of them is factually mistaken, four are predicted by the refined criterion, and two force a narrowing
of scope from heterotrophy to phagotrophy that the argument should have carried from the start.
\end{abstract}

\section{The observation}
\label{sec:obs}

The reason for eating is to gain access to the token supply held inside the agent being consumed.

There is a distinct evolutionary advantage to minimising the time spent \emph{exclusively} on energy
extraction. If some portion of the extraction can be carried on internally to the heterotroph --- so
that the heterotroph can meanwhile mate, seek further food, form hypotheses about its environment,
make predictions --- then extraction runs concurrently with everything else the agent does, and that
concurrency is worth paying for.

Hence eating has two phases:
\begin{enumerate}
\item the gross phase: getting the relevant material across the boundary and inside;
\item the internal phase: the remaining processes working on that material while the agent is engaged
elsewhere.
\end{enumerate}

The residue of the internal phase is the part that was not energy. It must leave. Material in, refuse
out: this makes the heterotroph, topologically, a tube.

The autotroph has no such need. Nothing about harvesting an external source calls for a tube
structure. And this is the arrangement found in the physical world.

\section{What the framework already supplies}
\label{sec:imports}

Every piece the argument needs is on the table. The claim of this note is not that a new mechanism is
required but that the tube \emph{follows}.

\paragraph{Access profiles.} The scientist note's trophic section distinguishes a source --- a
persistent emitter, rate-limited, non-exclusive, non-adaptive --- from prey --- a linear send,
exhaustible, exclusive, adaptive. Same sort, same rewrite rule, different multiplicity. Trophic type
is a fact about the cut, not a new kind of thing.

\paragraph{The asymmetry that generates the refuse.} A source emits \emph{tokens}. Prey is a
\emph{computation}: a structured term that happens to hold a stack. Harvesting a source yields token
and nothing else; breaking prey open yields token entangled with code.

\begin{observation}[Refuse is not an extra postulate]
\label{obs:refuse}
The residue exists precisely because prey are computations and sources are not. It is the non-token
part of a structured prey term. No assumption beyond Table~1 of the scientist note is required to
produce it.
\end{observation}

\paragraph{Media and interposed channels.} The engine note's error model generalises
\lstinline@for(y <- x)P | x!(Q)@ to \lstinline@for(y <- x1)P | C(x1,x2) | x2!(Q)@, and ranks media by
behaviour: relay, lossy, corrupting, reading, converting, with the engine at the top rung. A $d$-hop
pathway is already constructed there.

\paragraph{Porosity and individuation.} An individual is $\Chn_{\mathrm{in}}$ closed,
$\Chn_{\mathrm{boundary}}$ open. Boundary channels require redex pathways to internal channels or they
are stranding sites. Integration depth $d$ costs reliability $p^{d}$; ownership is what buys depth.

\paragraph{The tower.} The outside of one level is the inside of the next.

\paragraph{Namespace separation as the guarantor of non-interference.} In the game note, threat offers
had to live on their own names: a line signalling on a square's move channel would rendezvous with
that square's receipt and play the move. Perception must not be able to act.

\section{The central identification}
\label{sec:ident}

\begin{observation}[The gut is an internalised engine]
\label{obs:gut}
The digestive apparatus is a converting medium held inside the individual's boundary namespace: a
chain of interposed channels running from an intake channel to an expulsion channel, with the whole
chain owned.
\end{observation}

The mycorrhizal network was the external instance of the engine profile. The gut is the same profile
drawn on the other side of the ownership line. This is not an analogy between two systems; it is one
access profile appearing at two positions relative to a boundary.

The observation that makes the topology non-metaphorical is the next one.

\begin{observation}[The lumen is outside]
\label{obs:lumen}
The interior of the tube is not inside the individual. It is boundary namespace: owned, but open.
\end{observation}

A tube is therefore the construction that surrounds a region of the outside with the inside. That is
the tower's generating step, applied for a specific purpose: to acquire boundary surface without
acquiring exposure. The biological commonplace that the gut lumen is topologically external is, on
this reading, the same fact as the framework's fractality --- not a coincidence to be noted but the
mechanism itself.

\section{Statements to attempt}
\label{sec:props}

Numbering is provisional and every item below is a candidate rather than a result.

\begin{proposition}[Separation is not itself yielding]
\label{prop:A}
In the harvest of a structured prey term, the rewrite steps that separate token from residue are
metered but token-neutral. They cost and do not pay.
\end{proposition}

\begin{proposition}[The serialisation penalty]
\label{prop:B}
If separation is sequential with the scientist's control loop, the scientist pays its burn rate over
the whole separation with no inquiry, no foraging and no mating during it.
\end{proposition}

The target here is to write the foraging inequality with a sequential fraction $\sigma$ and derive the
Amdahl-shaped bound on sustainable harvest rate, so that the advantage of concurrency emerges as a
strict inequality with an explicit threshold rather than as a preference.

\begin{proposition}[Concurrency forces internalisation]
\label{prop:C}
Concurrency as such is trivially available in \rhoc{}: it is parallel composition. The content is
therefore not that extraction \emph{can} be concurrent but \emph{where the extraction channels live}.
If separation runs on boundary channels, other agents may rendezvous with partially processed
material. Exclusivity of a partially extracted harvest requires the separation chain to be owned.
\end{proposition}

Proposition~\ref{prop:C} is the ownership-buys-depth corollary specialised to digestion, and it is the
step that converts a remark about scheduling into a claim about topology.

\begin{proposition}[Two ends, not one]
\label{prop:D}
Intake and expulsion must sit on distinct names. On a single shared channel an offer of refuse and a
receipt for intake rendezvous: the organism re-ingests its residue, or the residue blocks the intake.
A single-opening organism is therefore not impossible but rate-limited by a mutual exclusion --- intake
and expulsion must alternate on the shared name.
\end{proposition}

This is the perception-must-not-act argument of the game note, one level up. It has an immediate
empirical check, and the check is favourable: cnidarians and flatworms possess a gastrovascular cavity
with a single opening and feed in discrete alternating bouts, while bilaterians with a through-gut feed
continuously. The framework should predict the throughput gap from the exclusion alone.

\begin{proposition}[Autotrophs are not tubes]
\label{prop:E}
For source access, both of the tube's ends lose their reason. The profile is non-exclusive, so nothing
is gained by fast acquisition; the source is unstructured, so there is no residue to expel.
\end{proposition}

The prediction is not merely that a plant needs no tube but that it should instead buy boundary area
directly. A leaf and a root system are that purchase.

\begin{corollary}[The second heterotroph's expense]
\label{cor:F}
The scientist note derives ``intelligence is a heterotroph's expense'' from the limit-of-inquiry
proposition together with structural recursion. The tube should be a second corollary of the same
proposition: exclusivity together with structure forces both purchases.
\end{corollary}

If Corollary~\ref{cor:F} goes through it is the strongest result available here, because it makes a
gross anatomical fact a consequence of an access profile. Intelligence and the tube become siblings
rather than neighbours.

\begin{proposition}[Transient versus persistent tubes]
\label{prop:G}
A food vacuole is a tube built and dissolved per meal; a gut is one maintained. Persistence carries a
maintenance cost against the metabolic stack, so there is a threshold harvest rate above which the
permanent tube beats the transient one.
\end{proposition}

Amoeba below the threshold, annelid above it. The composite-resolution proposition already says the
overlap region between strategies is populated by composites rather than atoms, so a graded answer here
is expected rather than a disappointment.

\begin{proposition}[Gut depth tracks incommensurability]
\label{prop:H}
Tokens held in a foreign colour are not nutritive as such; they must be converted. The number of
interposed converting stages should scale with the colour distance between prey namespace and consumer
namespace.
\end{proposition}

Carnivores eat close relatives and have short guts. Herbivores eat far ones and have long guts, often
renting whole foreign namespaces --- symbionts --- inside the boundary to perform conversions they
cannot perform themselves. A rented engine inside the tube is a market inside an individual, and it
survives the Coase objection recorded in the engine note precisely because, by
Observation~\ref{obs:lumen}, the lumen is not the inside.

\section{Rate of exchange: why breathing is not eating}
\label{sec:rate}

\emph{Dictated 11 September 2026, as a refinement of Section~\ref{sec:obs}.}

Respiration is also an exchange across a boundary. Something is admitted, a fraction of it is
extracted, and the remains are expelled. Yet the mammalian lung is a bag and not a tube. The dictated
intuition is that the discriminating variable is the \emph{rate} of exchange: digestive processes run
orders of magnitude slower than the respiratory cycle, and it is the slowness that forces the second
opening.

The intuition is pointing at something real. Rate as such is nevertheless the wrong variable, and
saying why turns out to sharpen the note rather than weaken it.

\subsection{Respiration goes tube exactly where one would not expect it}

Ventilation is unidirectional in fish --- water enters at the mouth and leaves at the opercular
openings --- and unidirectional in birds, where air moves the same way through the parabronchi during
both phases of the cycle. Unidirectional pulmonary flow is not confined to birds: it has been
demonstrated in alligators \cite{farmersanders2010}, in the savannah monitor \cite{schachner2014}, and
in iguanas, and Farmer's review argues from that distribution that the trait is ancestral for
diapsids and therefore \emph{not} an adaptation for high rates of gas exchange, since the animals
that have it include ectotherms that do not fly; candidate functions offered instead include reducing
the work of breathing, evaporative water loss and heat loss \cite{farmer2015}.

So a respiratory system can be a tube while running fast, and can be a bag while running fast. High
rate is neither necessary nor sufficient. What varies across those cases is the ratio of demand to
what a tidal exchanger can supply --- water holds a small fraction of air's oxygen and is far more
viscous --- rather than the rate itself.

\subsection{Two parameters, not one}

\begin{definition}[Occupancy ratio]
\label{def:lambda}
$\Lambda = \tau_{\mathrm{proc}} / \tau_{\mathrm{enc}}$, where $\tau_{\mathrm{proc}}$ is the residence
time of an admitted load inside the exchanger and $\tau_{\mathrm{enc}}$ is the interval at which a
further load could profitably be admitted.
\end{definition}

\begin{definition}[Separability]
\label{def:mu}
$\mu \in [0,1]$ measures the degree to which residue can leave past incoming material at the same
aperture without bulk transport: $\mu = 1$ for a residue miscible with the intake and free to leave by
diffusion, $\mu = 0$ for a residue that must be conveyed as a bolus.
\end{definition}

$\Lambda$ is what sets the sequential fraction $\sigma$ of Proposition~\ref{prop:B}: if the control
loop is blocked for the whole residence, $\sigma = \Lambda / (1 + \Lambda)$, and the Amdahl-shaped
bound bites only as $\Lambda$ approaches and passes unity. This is the sense in which the dictated
intuition is right. Slowness matters, but only relative to the opportunity structure. A process that
takes a day is cheap to serialise if the next load cannot arrive for a week.

\subsection{Breathing fails on separability, not on rate}

The respiratory residue fraction is close to one: almost all of the mass admitted is expelled again.
Bag topology survives that because $\mu \approx 1$. The mammalian lung in fact tolerates gross
re-mixing --- an inspired bolus on the order of a quarter of a litre mixes with something like a litre
and a half of residual gas, six parts spent medium to one part fresh, every cycle
\cite{farmer2015} --- and it works, because the residue is the same phase as the resource and
separation happens by diffusion at the membrane rather than by transport through the aperture.

Digestion has $\mu \approx 0$. A bolus of partly extracted material cannot diffuse past incoming food,
and the cost of re-mixing is not a few percent of exchange efficiency but contamination of the batch.

\begin{proposition}[The tube criterion]
\label{prop:I}
A tube is forced when the agent is \emph{phagotrophic} --- admitting bulk material with a
non-absorbable fraction --- and $\Lambda \gtrsim 1$ and $\mu \approx 0$. Failure of any one of the
three is sufficient for a bag.
\end{proposition}

Proposition~\ref{prop:I} subsumes Proposition~\ref{prop:D} rather than replacing it: the mutual
exclusion on a shared name is what the agent pays, and $\Lambda$ is how much of the cycle that
exclusion consumes.

\begin{corollary}[Amendment to Proposition~\ref{prop:E}]
\label{cor:E2}
Proposition~\ref{prop:E} was stated for autotrophs and should have been stated for non-phagotrophs.
Source access is one way to have nothing to expel; it is not the only way.
\end{corollary}

\section{Counterexamples, and where they land}
\label{sec:stay}

M.\ Stay (personal communication, 11 September 2026) proposed a list of animals with bag-like rather
than tube-like digestive systems --- cnidarians, ctenophores, flatworms, xenacoelomorphs and
ophiuroids --- together with two cases from outside the animals: fungi, which are heterotrophs, and
plants such as \emph{Monotropa uniflora}, which parasitise fungi rather than photosynthesise. The list
is the right test set, and it separates cleanly into three kinds of case: one that is factually
mistaken, four that the occupancy ratio predicts, and two that force a correction of scope.

\begin{center}
\begin{tabular}{@{}llp{5.2cm}@{}}
\toprule
Case & Verdict & Why \\
\midrule
Ctenophora & fails & has a through-gut \\
Cnidaria & predicted & low $\Lambda$; duty-cycles when $\Lambda$ rises \\
Platyhelminthes & predicted & buys area; gutless where osmotrophic \\
Xenacoelomorpha & predicted & no lumen at all; diffusion-scale bodies \\
Ophiuroidea & predicted & secondary loss, microphagous diets \\
Fungi & scope & osmotroph, not phagotroph \\
Mycoheterotrophs & scope & osmotroph, not phagotroph \\
\bottomrule
\end{tabular}
\end{center}

\subsection{Ctenophora: the counterexample that fails}

Comb jellies are not bags. Presnell and colleagues showed by time-lapse imaging that the ctenophore
gut is unidirectional and functionally tripartite, with waste expelled through terminal anal pores
that are specialised to control outflow, resolving a long-standing misreading of those pores that had
supported the blind-gut picture \cite{presnell2016,giribet2016}.

This is more than the removal of a counterexample. Ctenophores branch very deep in the metazoan tree,
so a through-gut there is evidence that the construction is not an inheritance of the bilaterian body
plan but something arrived at from an access profile --- which is the note's thesis.

\subsection{Cnidaria: the predicted duty cycle}

Mean digestion time in \emph{Aurelia aurita} has been measured at about one hour
\cite{ishii2001}. Against the encounter interval of a drifting tentacle feeder on dilute plankton this
is a low $\Lambda$, and Proposition~\ref{prop:I} accordingly predicts a bag.

The prediction with teeth is what should happen when $\Lambda$ rises, and it has been observed. In
\emph{Pelagia noctiluca} ephyrae feeding on a dense and pulsed prey field, the gut saturates in about
fifteen minutes while digestion takes about eighteen hours, and the authors read the mismatch as
implying a diel feeding periodicity \cite{gordoa2013}. Fast fill, slow clear, forced alternation: that
is Proposition~\ref{prop:D} observed in a single organism, with the ratio measured.

\subsection{Platyhelminthes and Xenacoelomorpha: area instead of length}

Acoels have no gut lumen at all. The mouth opens into a syncytial digestive parenchyma without
epithelial lining, in which particles are handled by phagocytosis. At body scales where the diffusion
distance is a few cell diameters there is no bulk transport to arrange and $\mu$ has no work to do.
The branched blind gut of a triclad is the same move at larger scale: it buys boundary area rather
than a second opening, which is the purchase Proposition~\ref{prop:E} assigns to autotrophs.

Cestodes are the sharper case. Tapeworms have no digestive system whatever and take up small organic
monomers across the tegument. A cestode is an animal, a bilaterian, a heterotroph, and descended from
ancestors with a through-gut --- and it is not a tube, because something else did the digestion. No
formulation in terms of heterotrophy can accommodate that; Proposition~\ref{prop:I} accommodates it
by the phagotrophy clause.

\subsection{Ophiuroidea: secondary loss under a diet shift}

Brittle stars lack an anus, as do paxillosid asteroids, and the review literature catalogues repeated
independent losses of the anal opening across Bilateria, of which these are two
\cite{hejnol2015,jangoux1982}. The direction of change matters: this is loss from a through-gut
ancestor rather than a lineage that never built one, and it is concentrated in microphagous habits,
suspension and deposit feeding on small particles, where the indigestible fraction per load is small
and $\Lambda$ correspondingly low.

That yields a within-clade test. Macrophagous ophiuroids --- scavengers and predators taking large
items --- should show markedly harder duty-cycling than their microphagous congeners, and if they do
not, Proposition~\ref{prop:I} is in trouble on ground of its own choosing.

\subsection{Fungi and mycoheterotrophs: the scope was stated too widely}

These are not counterexamples. They are a correction, and a useful one, because they identify the
clause the note had left implicit.

Heterotrophy is a claim about carbon source. The tube argument is a claim about \emph{phagotrophy}:
about admitting bulk material that contains something which will not be absorbed. A fungus secretes
its enzymes outward and takes up monomers; the residue never crosses the boundary at all.
Mycoheterotrophic plants do the same at one further remove, drawing fixed carbon through a
mycorrhizal interface \cite{leake1994,bidartondo2005}, and \emph{Monotropa uniflora} is the
resolutive case, taking everything it has through associations with Russulaceae.

\begin{observation}[Osmotrophs evert the lumen]
\label{obs:evert}
Observation~\ref{obs:lumen} said the lumen is boundary namespace, owned but open. An osmotroph runs
the same converting chain without the ownership: the substrate \emph{is} the lumen, and the engine
sits at the boundary rather than folded inside it. The fungus is the tube construction turned inside
out --- which is why the mycorrhizal network appeared in the engine note as the external instance of
the very profile Observation~\ref{obs:gut} draws internally.
\end{observation}

The cost structure is then legible rather than anomalous. Internalising the chain buys exclusivity
over a partially extracted harvest, which is Proposition~\ref{prop:C}; externalising it surrenders
exclusivity and in exchange never has to transport residue. An osmotroph is an agent that has taken
the other side of that trade, and it should be predicted to be vulnerable exactly where the trade is
worst, namely to theft of extracellular product by neighbours.

\subsection{A prior statement of Proposition~\ref{prop:B}}

Honesty about priority: the functional claim is already in the literature. The same review that
catalogues the losses states in passing that a one-way gut processes food more efficiently and permits
uptake while the animal is still digesting \cite{hejnol2015}. The contribution claimed here is
therefore not the observation but the derivation --- that the two openings follow from concurrency
together with ownership, and that the threshold is a computable function of $\Lambda$ and $\mu$ rather
than a qualitative preference.

\section{What to write down formally}
\label{sec:formal}

\paragraph{The two-phase harvest as a term.} An exclusive acquisition rendezvous on a boundary
channel, in parallel with a persistent digestion chain on internal channels, terminating in a send on
the expulsion channel. The existing \texttt{Forage} of the scientist note's worked ecosystem and the
\texttt{Chain} of \texttt{engine.rho} supply most of it.

\paragraph{The tube as a namespace shape.} $\Chn_{\mathrm{boundary}}$ factors as
$\Chn_{\mathrm{in}} \sqcup \Chn_{\mathrm{out}}$ with a directed redex pathway from in to out and no
return path. Then \emph{tube} is a formula in the logic rather than a designation, which is the same
move already made for ``being a scientist is a depth-$\geq 2$ dialogue loop through a both-type
channel''.

\paragraph{Whether the absence of a return path is forced.} Coprophagy, caecotrophy and rumination are
return paths, and they occur exactly where one pass leaves too much unextracted. This is cheap to state
as a condition on residual yield and should be stated, since a proposition that forbids observed
behaviour is worse than one that predicts it.

\paragraph{The occupancy ratio as a term-level quantity.} $\Lambda$ of
Definition~\ref{def:lambda} should be recoverable from the harvest term itself: $\tau_{\mathrm{proc}}$
as the metered cost of the digestion chain, $\tau_{\mathrm{enc}}$ as the expected waiting time on the
intake channel under the ambient offer rate. Then Proposition~\ref{prop:I} becomes an inequality
between two quantities the framework already meters, and the tube threshold is derived rather than
posited.

\paragraph{Separability as a namespace condition.} $\mu$ is not a physical parameter but a statement
about whether residue and intake can share a name without rendezvous --- that is, whether the sorts
involved make the crossed redex impossible. Gas exchange has $\mu = 1$ because the spent medium and
the resource are the same sort and the extraction happens at the wall rather than at the aperture.
This wants stating as a condition on the sort discipline, at which point Proposition~\ref{prop:I} is
entirely internal.

\paragraph{The physical topology as a shadow.} A through-gut makes the body a genus-one surface. That
is real, but it is not the content. The content is the channel topology; the surface is what an
embodiment of that topology in three-space looks like. The note should say so explicitly, so that it
does not appear to be arguing from anatomy to mathematics.

\section{Open questions}
\label{sec:open}

\begin{question}
Is the tube forced or merely favoured? Composite resolution predicts mixotrophs, and predicts that
their tubes should be facultative. Is that borne out?
\end{question}

\begin{question}
Does the surface-area argument connect to the critical radius $r^{*}$ of the compression-versus-autonomy
result? A tube buys boundary area at fixed volume; $r^{*}$ is a critical radius. There is plausibly one
calculation here rather than two.
\end{question}

\begin{question}
Where does circulation enter? Beyond a certain body size the tube's yield must be distributed, which is
a second internal medium with a different profile --- relay rather than converting --- and a different
topology, branching and returning. If the tube follows from the access profile, does the vascular tree
follow from the tube together with the allometric constraint?
\end{question}

\begin{question}
What is the assay analogue? Tasting is an assay run at the intake channel deciding whether to admit;
vomiting is a refutation acted upon after admission. The trichotomy and budget-relative refutation
should both apply, and the lossy-medium correction owed upstream to the scientist note's assay section
may turn out to be the same correction.
\end{question}

\begin{question}
Does the tube give a new grading dimension, or is it already covered by the vector over typed
namespaces? Suspicion: the tube is a fact about $k_{\mathrm{chn}}$, which would be an independent
argument for the fifth namespace.
\end{question}

\begin{question}
Does the osmotroph/phagotroph split have a framework-native characterisation? Observation~\ref{obs:evert}
says the osmotroph runs the converting chain outside the ownership boundary. If that is right, the
split should fall out of where the chain sits relative to $\Chn_{\mathrm{boundary}}$, and osmotrophy
should carry a predictable exposure to theft that phagotrophy does not.
\end{question}

\begin{question}
Is there a second critical point at high $\Lambda$? Ruminants and hindgut fermenters run
$\tau_{\mathrm{proc}}$ far above $\tau_{\mathrm{enc}}$ and answer with buffering --- a sequence of
chambers, and in the ruminant a deliberate return path. Buffering is the standard answer to a
throughput mismatch, and the question is whether the framework predicts the chamber count.
\end{question}

\section{Filing}
\label{sec:filing}

The natural home is the trophic-structure section of the scientist note, since the material is short
and is a corollary of what is there. If Propositions~\ref{prop:B}, \ref{prop:D} and \ref{prop:H} develop
enough measurable content --- a threshold, a throughput ratio, a scaling law --- the note stands on its
own in the rho-life series instead.


\begin{thebibliography}{99}

\bibitem{bidartondo2005} Bidartondo, M.~I. (2005). The evolutionary ecology of myco-heterotrophy.
\emph{New Phytologist} 167: 335--352.

\bibitem{farmer2015} Farmer, C.~G. (2015). The evolution of unidirectional pulmonary airflow.
\emph{Physiology} 30. \texttt{doi:10.1152/physiol.00056.2014}

\bibitem{farmersanders2010} Farmer, C.~G. and Sanders, K. (2010). Unidirectional airflow in the lungs
of alligators. \emph{Science} 327: 338--340. \texttt{doi:10.1126/science.1180219}

\bibitem{giribet2016} Giribet, G. (2016). Zoology: at last an exit for ctenophores.
\emph{Current Biology} 26: R918--R920. \texttt{doi:10.1016/j.cub.2016.08.036}

\bibitem{gordoa2013} Gordoa, A., Acu\~na, J.~L., Farr\'es, R. and Bacher, K. (2013). Burst feeding of
\emph{Pelagia noctiluca} ephyrae on Atlantic bluefin tuna (\emph{Thunnus thynnus}) eggs.
\emph{PLoS ONE} 8(9): e74721. \texttt{doi:10.1371/journal.pone.0074721}

\bibitem{hejnol2015} Hejnol, A. and Mart\'in-Dur\'an, J.~M. (2015). Getting to the bottom of anal
evolution. \emph{Zoologischer Anzeiger} 256: 61--74. \texttt{doi:10.1016/j.jcz.2015.02.006}

\bibitem{ishii2001} Ishii, H. and Tanaka, F. (2001). Food and feeding of \emph{Aurelia aurita} in
Tokyo Bay with an analysis of stomach contents and a measurement of digestion times.
\emph{Hydrobiologia} 451: 311--320.

\bibitem{jangoux1982} Jangoux, M. (1982). Digestive systems: Ophiuroidea. In Jangoux, M. and
Lawrence, J.~M. (eds.), \emph{Echinoderm Nutrition}. Balkema, Rotterdam.

\bibitem{leake1994} Leake, J.~R. (1994). The biology of myco-heterotrophic (`saprophytic') plants.
\emph{New Phytologist} 127: 171--216.

\bibitem{presnell2016} Presnell, J.~S., Vandepas, L.~E., Warren, K.~J., Swalla, B.~J., Amemiya, C.~T.
and Browne, W.~E. (2016). The presence of a functionally tripartite through-gut in Ctenophora has
implications for metazoan character trait evolution. \emph{Current Biology} 26: 2814--2820.
\texttt{doi:10.1016/j.cub.2016.08.019}

\bibitem{schachner2014} Schachner, E.~R., Cieri, R.~L., Butler, J.~P. and Farmer, C.~G. (2014).
Unidirectional pulmonary airflow patterns in the savannah monitor lizard. \emph{Nature} 506: 367--370.

\end{thebibliography}

\end{document}
